\documentclass[aspectratio=169,10pt]{beamer}
\usetheme{metropolis}
\usepackage{booktabs,amsmath,amssymb}
\title{Q30: Inside PyTorch}
\subtitle{torch.compile layers + autodiff memory}
\author{DLS 2026 --- Kukla}
\date{October 2026}
\begin{document}
\maketitle

\begin{frame}{The question}
\begin{block}{Student question (\#29)}
How does the PyTorch compiler work and what are its layers? Why does functional autodiff use far more
memory than tape-based autodiff?
\end{block}
\end{frame}
\begin{frame}{Eager vs.\ compiled}
\begin{itemize}
\item Default: \textbf{eager} --- ops dispatch immediately, easy to debug, slow.
\item \texttt{torch.compile}: JIT-compiles, keeps eager semantics, recovers speed.
\item Philosophy: \emph{capture what you can, fall back where you must}.
\end{itemize}
\end{frame}
\begin{frame}{The torch.compile stack}
\begin{itemize}
\item \textbf{TorchDynamo}: trace bytecode $\to$ FX graph; \emph{graph breaks} on untraceable code.
\item \textbf{AOTAutograd}: capture forward \emph{and} backward ahead of time.
\item \textbf{TorchInductor}: fuse ops, lower to kernels.
\item \textbf{Triton}: the emitted GPU kernel language.
\end{itemize}
\end{frame}
\begin{frame}{Two flavors of autodiff}
\begin{itemize}
\item \textbf{Reverse-mode / tape}: record a tape forward, walk it backward (classic autograd).
\item \textbf{Functional}: transform a pure fn into its gradient fn (\texttt{grad}, \texttt{vmap}, \texttt{jit}).
\end{itemize}
\end{frame}
\begin{frame}{Why functional AD can use more memory}
\begin{itemize}
\item \textbf{Activations} saved for backward dominate training memory.
\item A naive functional form may retain more intermediates than a tuned tape.
\item Fix: \textbf{gradient checkpointing} --- save few activations, \emph{recompute} the rest.
\end{itemize}
\alert{Trade extra compute for a big memory cut (can turn OOM into trainable).}
\end{frame}

\end{document}
